\scanpage{149}
On suppose donné trois catégories \inserted{additives} $\mathscr V,\mathscr V',\mathscr V''$, d'où $\mathcal M,\mathcal M',\mathcal M''$, on suppose donné un foncteur biadditif $\mathscr V\times\mathscr V'\unsure{\mathord{-}|\longrightarrow}\mathscr V''$. Alors il se prolonge de façon ess\textsuperscript{t} unique en un foncteur bi-additif $\mathcal M\times\mathcal M'-\mathcal M''$. En fait, \unsure{on|ce} voit que la propriété \inserted{pseudo-}universelle qui définit $\mathcal M$ en terme de $\mathscr V$, s'étend aux bifoncteurs, ......

On \deleted{\unsure{y}} utilisera cette construction surtout pour un bifoncteur $\mathscr V\times\mathscr V-\mathscr V$ ``\uline{symétrique}'', et ``\uline{associatif}''\inserted{ et ``unitaire''}, noté $\deleted{\otimes}\boxtimes$. Alors on trouve un bifoncteur de $=$ nature dans $\mathcal M$, noté $\otimes$. Si $k$ est \uline{de car. $0$}\inserted{ (p.ex. $\mathbb Q$)}, on peut construire alors toutes les opérations tensorielles\inserted{ covariantes} dans $\mathcal M$ [en plusieurs arguments], par exemple on pourra définir les
\[
 \overset{\unsure{n|-}}{\bigotimes} M,\quad \overset{n}{\bigwedge}M,\quad\deleted{\text{etc}}\ S^n(M)\quad\text{etc ....}
\]
% archive sheet number: 148
